\chapter{实验评估}

\section{Evaluation Goals}
本章实验回答三个问题。第一，FTCL 的 language runtime 是否保持了预期语义？第二，CPU 和 CUDA geometry backends 是否能产生等价结果？第三，从 FTCL script 用户的角度看，CUDA 在什么规模下能够带来 end-to-end speedup？

项目的实验数据、CSV 文件和图表生成脚本都位于 \texttt{docs} 目录下。测量内容包括 semantic pass rate、parser backend timing、thread-channel latency、frame-time distribution、CPU/CUDA geometry speedup、throughput、numerical error、break-even point、ablation data 和 concurrent geometry performance。这些指标共同覆盖 correctness、performance 和 scalability 三个层面。

\section{Experimental Methodology}
除非特别说明，本文所有 performance results 都被视为 end-to-end system measurements。一次 geometry command measurement 包含 FTCL command dispatch、argument parsing、UVec handle lookup、device synchronization、output UVec allocation、backend execution，以及 CUDA kernel launch 后的 synchronization。它不是纯 CUDA kernel time。本文有意采用这种测量方式，因为论文评价的是脚本用户实际看到的系统性能。

Benchmark pipeline 遵循以下原则：

\begin{enumerate}
  \item \textbf{使用 Release builds。} Debug builds 不进入性能讨论，因为 interpreter dispatch 和 STL containers 对 optimization level 非常敏感。
  \item \textbf{区分 warmup 和 measurement。} Parser backend timing 和 geometry timing scripts 在可支持的情况下会先执行 warmup runs，再记录 samples。
  \item \textbf{优先使用 median 或 percentile。} Latency distributions 使用 P50、P95 和 P99 描述；geometry timing 在有多个样本时使用 representative median command latency。
  \item \textbf{单独看待 readback cost。} 部分实验将结果保留在 UVec handles 中，ablation experiments 则显式测量将 large UVec outputs 转回 script lists 的成本。
  \item \textbf{先比较正确性，再解释速度。} 只有 CPU/GPU equivalence tests 在数值容忍范围内通过时，CUDA speedups 才有意义。
\end{enumerate}

主要派生指标定义如下：
\[
  \mathrm{speedup} = \frac{T_{\mathrm{CPU}}}{T_{\mathrm{CUDA}}},
  \quad
  \mathrm{throughput} = \frac{\mathrm{work\ items}}{T},
  \quad
  \mathrm{work\ items} = N \times Q.
\]
对于 distance matrix、nearest point 和 range-count workloads，\(N\) 表示 data points 数量，\(Q\) 表示 query points 或 query regions 数量。

\subsection{CPU Baseline Definition}
本文的 CPU baseline 是 FTCL CPU geometry backend，而不是手工优化的 parallel CPU library。这个区别很重要：实验回答的问题是脚本用户会问的问题，即同一个 FTCL command 在 CPU device 和 CUDA device 上执行时，end-to-end improvement 是多少。

因此，CPU row 和 CUDA row 尽量共享相同的 script-level command path。主性能实验中，输入会在计时前预先构造为 UVec handles；计时区间包括 command dispatch、handle lookup、output UVec allocation、backend execution 和必要 synchronization。CPU row 使用 C++ scalar loops 处理 CPU-resident UVec data；CUDA row 使用 CUDA device buffers 和 kernels，并且当前 CUDA helper 在 kernel launch 后调用 \texttt{cudaDeviceSynchronize}，保证 command 返回时 backend work 已完成。

这个 baseline 对评价 FTCL 作为 heterogeneous scripting system 是公平的，但不能被解释为对抗高度优化的 all-core CPU geometry library。实验服务器拥有 128 个 logical CPU threads，而当前 CPU geometry backend 故意保持简单。因此，报告的 speedup 衡量的是在相同 command interface 下，将 FTCL backend 从 scalar CPU implementation 移动到 CUDA implementation 后获得的收益。未来加入 all-core CPU backend 会是有价值的补充 baseline。

\begin{table}[H]
\centering
\caption{实验环境信息（Experimental environment for reproducibility）}
\label{tab:environment}
\begin{tabularx}{\textwidth}{lX}
\toprule
Item & Recorded value \\
\midrule
Operating system & WSL Ubuntu environment used for build, tests, CUDA execution, and figure generation \\
CPU model & Intel Xeon Platinum 8358 CPU at 2.60 GHz \\
CPU topology & 2 sockets, 32 cores per socket, 2 hardware threads per core, 128 logical CPUs \\
CPU frequency & 800 MHz minimum, 3.40 GHz maximum as reported by \texttt{lscpu} \\
NUMA layout & 2 NUMA nodes; node 0 contains CPUs 0--31 and 64--95, node 1 contains CPUs 32--63 and 96--127 \\
Main memory & 1.0 TiB total memory; about 623 GiB available when the environment was recorded \\
GPU model & 8 NVIDIA A800-SXM4-80GB GPUs \\
GPU memory & 81920 MiB per GPU, approximately 640 GiB aggregate device memory \\
GPU driver & NVIDIA driver 530.30.02 \\
CUDA compute capability & 8.0 for each A800 GPU \\
Build system & CMake with Release configuration \\
C++ standard & C++20, selected by \texttt{FTCL\_CXX\_STANDARD=20} \\
Compiler & GCC or Clang available inside WSL \\
CUDA option & \texttt{FTCL\_ENABLE\_CUDA=ON} for CUDA geometry experiments \\
CPU build path & \texttt{build} or \texttt{/tmp/ftcl-build-geometry-cpu} \\
CUDA build path & \texttt{/tmp/ftcl-build-geometry-cuda} \\
Benchmark data & CSV files under \texttt{docs/data/benchmarks} and \texttt{docs/data/geometry} \\
Generated figures & SVG/PNG figures under \texttt{docs/figures} and \texttt{contents/figures} \\
Important note & Free GPU memory is time-dependent because other users or jobs may share the same server \\
\bottomrule
\end{tabularx}
\end{table}

表 \ref{tab:environment} 是环境记录，而不是硬件无关结论。本文所有 performance numbers 都应理解为在该 multi-GPU server 上的实测结果。服务器的 CPU core count 也说明 CPU baseline 并不来自弱硬件；真正的实验边界在于当前 CPU backend 采用 scalar implementation，而不是 fully parallel CPU implementation。

\section{Semantic Regression Tests}
Semantic regression tests 覆盖 core commands、collections、control flow、expressions 和 threading。图 \ref{fig:semantic-pass} 显示当前测量的 test suites 达到 100 percent pass rate。

\begin{figure}[H]
  \centering
  \includegraphics[width=0.78\textwidth]{semantic-pass-rate.png}
  \caption{测试套件语义通过率（Semantic pass rate by test suite）}
  \label{fig:semantic-pass}
\end{figure}

这个结果并不表示 FTCL 已经完整实现 Tcl。它表示本文选定的 Tcl subset 已经形成稳定 regression baseline。对于 parser 和 expression 这类容易引入细微语义变化的模块，这一点尤其重要。

\section{Parser Backend Timing}
图 \ref{fig:parser-backend} 比较 legacy parser 和 token-stream parser 在 Tcl subset benchmark 上的耗时。Token-stream path 并不一定更快，因为它增加了 lexing pass，并保存 token metadata。它的价值主要是架构层面的：更好的 diagnostics、更清晰的 command boundaries，以及通向 AST 和 bytecode execution 的路径。

\begin{figure}[H]
  \centering
  \includegraphics[width=0.80\textwidth]{parser-backend-timing.png}
  \caption{Parser backend 耗时对比（Timing comparison between parser backends）}
  \label{fig:parser-backend}
\end{figure}

直接结论：token-stream parser 的短期意义不是速度，而是可维护性、错误定位和后续编译路径。

\section{Concurrency and Runtime Smoothness}
Thread-channel latency 和 frame-time distribution 如图 \ref{fig:channel-latency} 与图 \ref{fig:frame-time} 所示。Channel experiment 通过 ping-pong round trip 估计 one-way latency；frame-time experiment 测量一个 non-interactive game-like workload。

\begin{figure}[H]
  \centering
  \includegraphics[width=0.78\textwidth]{channel-latency.png}
  \caption{线程通道单向延迟分布（Thread-channel one-way latency distribution）}
  \label{fig:channel-latency}
\end{figure}

\begin{figure}[H]
  \centering
  \includegraphics[width=0.78\textwidth]{frame-time.png}
  \caption{类游戏 workload 的帧耗时分布（Frame-time distribution）}
  \label{fig:frame-time}
\end{figure}

这些实验说明 FTCL 可以支持简单 concurrent scripts 和 interactive demos。相比平均值，P50、P95 和 P99 这类 tail-latency metrics 更能反映脚本系统的实际交互体验。

\section{CPU/GPU Equivalence}
对 CUDA backend 而言，correctness 比 speed 更重要。Equivalence tests 会让同一 geometry workload 同时在 CPU 和 CUDA 上运行，比较返回的 UVec values，并允许较小 floating-point tolerance。图 \ref{fig:error-dist} 展示 numerical error distribution。

\begin{figure}[H]
  \centering
  \includegraphics[width=0.70\textwidth]{geometry-error-distribution.png}
  \caption{CPU 与 CUDA 几何结果的数值误差分布（Numerical error distribution）}
  \label{fig:error-dist}
\end{figure}

生成数据中的 observed errors 接近零。这支持了本文的基本正确性主张：在测试覆盖到的 workloads 上，CUDA implementation 保持了 CPU geometry semantics。

\section{Geometry Performance Scaling}
性能实验聚焦三个代表性算法：distance matrix、nearest point 和 circle range count。它们的共同点是存在大量 independent work items。表 \ref{tab:geom-speedup} 根据保存的 performance CSV 总结 end-to-end speedup。

Single-GPU performance-scale experiment 使用一个 FTCL interpreter 和一个 CUDA device。每个 scale 下，benchmark 会在计时前构造 CPU 和 CUDA UVec inputs，执行可选 untimed warmup，并报告 median command time。这意味着 input point-list parsing cost 被排除在主计时之外，而 output allocation 和 backend synchronization 仍然包含在内。实验目标是测量 prebuilt UVec handles 在 realistic script pipeline 中的使用成本。

\begin{table}[H]
\centering
\caption{代表性几何算法的 end-to-end CPU/CUDA speedup}
\label{tab:geom-speedup}
\begin{tabular}{lrrr}
\toprule
Algorithm & N,Q & Work items & Speedup \\
\midrule
Distance matrix & 2048, 256 & 524288 & 6.77x \\
Nearest point & 2048, 256 & 524288 & 2.05x \\
Circle range count & 2048, 256 & 524288 & 1.41x \\
Distance matrix & 8192, 1024 & 8388608 & 51.90x \\
Nearest point & 8192, 1024 & 8388608 & 8.39x \\
Circle range count & 8192, 1024 & 8388608 & 6.50x \\
Distance matrix & 32768, 2048 & 67108864 & 224.00x \\
Nearest point & 32768, 2048 & 67108864 & 16.55x \\
Circle range count & 32768, 2048 & 67108864 & 13.13x \\
\bottomrule
\end{tabular}
\end{table}

\begin{figure}[H]
  \centering
  \includegraphics[width=0.82\textwidth]{geometry-speedup.png}
  \caption{批量几何算法的 CPU/CUDA speedup}
  \label{fig:geom-speedup}
\end{figure}

图 \ref{fig:execution-bars} 和图 \ref{fig:throughput} 展示 execution time 和 throughput。Distance matrix 获得最高 speedup，因为它的 work items 简单且彼此独立。Nearest point 和 range count 同样受益于 CUDA，但每个 query 需要更多 memory reads，因此 scaling pattern 不完全相同。

\begin{figure}[H]
  \centering
  \includegraphics[width=0.82\textwidth]{geometry-execution-bars.png}
  \caption{最大测量规模下的执行时间对比（Execution-time comparison）}
  \label{fig:execution-bars}
\end{figure}

\begin{figure}[H]
  \centering
  \includegraphics[width=0.82\textwidth]{geometry-throughput.png}
  \caption{几何算法吞吐量曲线（Throughput curve）}
  \label{fig:throughput}
\end{figure}

直接结论：work items 足够多且输出紧凑时，CUDA backend 的优势更稳定；输出规模过大时，allocation 和 memory bandwidth 会成为瓶颈。

\section{Multi-GPU Scaling}
实验服务器包含 8 张 A800 GPUs，因此本文进一步测量 FTCL 是否真的能把多个 CUDA devices 暴露给 script-level geometry workloads。Multi-GPU benchmark 使用 weak scaling：每张 GPU 接收相同数量的 query points，而总 query count 随 GPU 数量增长。每个 worker 拥有独立 FTCL interpreter，在不同 CUDA device 上创建 UVec inputs，并分别在 \texttt{cuda:0}、\texttt{cuda:1}、\ldots、\texttt{cuda:7} 上启动相同 geometry command。报告的 speedup 是相对于 1 GPU 的 throughput speedup。

Benchmark 运行时需要让所有 devices 可见，例如：

\begin{lstlisting}[language=bash]
CUDA_VISIBLE_DEVICES=0,1,2,3,4,5,6,7 \
FTCL_MULTI_GPU_MODE=paper \
FTCL_MULTI_GPU_SCALING=weak \
test/bench_geometry_multi_gpu
\end{lstlisting}

记录运行中，\(N=32768\)，每张 GPU 接收 \(Q=4096\) 个 query objects，每行报告三轮 median。因此 1、2、4、8 GPUs 对应的 total query count 分别是 \(4096\)、\(8192\)、\(16384\) 和 \(32768\)。Timer 测量从同时释放所有 worker threads 到所有 workers 获得 output UVec handles 的 wall-clock interval。Length checking 和 output destruction 在计时区间之后执行，因此计时聚焦 command dispatch、UVec handle lookup、output allocation、CUDA launch、backend synchronization 和 result-handle creation。

这个实验还暴露了一个重要 runtime bottleneck。原始 UVec command manager 使用一个 global mutex 保护整个 handle table，并在持锁状态下执行 geometry commands。这会导致彼此独立的 CUDA commands 在 script runtime layer 被串行化。实现后来改为：global mutex 只保护 handle-table lookup，每个 UVec handle 再拥有独立 entry mutex。这样，不同 GPU 上的独立 UVec handles 可以并发执行。

\begin{table}[H]
\centering
\caption{多 GPU 几何命令 weak-scaling speedup}
\label{tab:multi-gpu-speedup}
\begin{tabular}{lrrrr}
\toprule
Algorithm & 1 GPU & 2 GPUs & 4 GPUs & 8 GPUs \\
\midrule
Distance matrix & 1.00x & 0.96x & 0.97x & 1.12x \\
Nearest point & 1.00x & 2.00x & 3.99x & 3.74x \\
Circle range count & 1.00x & 1.98x & 3.97x & 7.90x \\
\bottomrule
\end{tabular}
\end{table}

\begin{figure}[H]
  \centering
  \includegraphics[width=0.82\textwidth]{geometry-multi-gpu-scaling.png}
  \caption{1、2、4、8 张 A800 GPU 的 weak-scaling speedup}
  \label{fig:multi-gpu-scaling}
\end{figure}

直接结论：circle range count 接近 linear multi-GPU throughput scaling，而 distance matrix 暴露出 output-dominant bottleneck。

结果说明 multi-GPU execution 有价值，但强烈依赖 algorithm shape。Circle range counting 几乎线性扩展到 8 GPUs，因为每个 query 只产生一个 compact result，并且 work pattern 规则。Nearest point 扩展到 4 GPUs 表现良好，但到 8 GPUs 后下降，说明 runtime overhead、allocation 和 synchronization 开始占据主导。Distance matrix 的 multi-GPU scaling 最弱，因为每个 query 会产生 large output row；output allocation 和 global memory traffic 主导了 command time。这是一个有价值的 negative result：系统已经可以使用 8 GPUs，但要让 large-output kernels 真正扩展，还需要 memory pools、persistent output buffers 和 CUDA streams。

\section{Break-even Analysis}
CUDA acceleration 存在 fixed overhead。小规模 workload 可能在 CUDA 上更慢，因为 kernel launch、synchronization、allocation 和 command dispatch 会压过真正的计算量。图 \ref{fig:break-even} 使用 work items \(N \times Q\) 作为横轴，speedup 作为纵轴。

\begin{figure}[H]
  \centering
  \includegraphics[width=0.82\textwidth]{geometry-break-even.png}
  \caption{按 work items 统计的 CUDA break-even analysis}
  \label{fig:break-even}
\end{figure}

直接结论：只有 work-item count 足以摊销 launch、allocation、synchronization 和 UVec management costs 后，CUDA 才真正值得使用。

保存的数据估计 distance matrix 的 break-even point 约为 8192 work items，nearest point 约为 47848，circle range count 约为 121209。这些值不是 universal constants。它们依赖 hardware、driver state、input size、output size 以及 measurement 是否包含 readback。

\section{Ablation and Concurrent Geometry}
Ablation experiment 将 inline literal parsing、prebuilt UVec handles 和 readback overhead 分离。实验显示，prebuilt handles 可以显著降低 command latency，尤其对 CUDA commands 更明显。同时，将 large outputs 转回 script list 可能支配总时间。因此，高性能 FTCL script 应尽量让数据在 UVec form 中跨多个 geometry commands 保持驻留，只在最后读取 summary results。

\begin{table}[H]
\centering
\caption{几何消融实验解释（Interpretation of geometry ablation experiment）}
\label{tab:ablation-interpretation}
\begin{tabularx}{\textwidth}{lX}
\toprule
Experiment mode & Interpretation \\
\midrule
Inline literal input & 测量从 script text 解析 point lists、创建临时数据结构、分发 command 并执行 backend 的完整成本。该模式方便，但大输入时容易被 script parsing 主导。 \\
Prebuilt UVec without readback & 测量 handle lookup、device synchronization、output allocation 和 backend execution，同时将数据保持在 UVec form 中。该模式最接近 high-throughput script pipeline。 \\
Prebuilt UVec with readback & 额外加入将 output UVec 转为 script list 的成本。对于 distance matrices 这类 large outputs，readback 会主导总时间并掩盖 CUDA computation 的收益。 \\
\bottomrule
\end{tabularx}
\end{table}

消融实验给出了一条实践原则：FTCL scripts 应把 UVec handles 当作 long-lived data objects 使用，避免把 large intermediate results 转回 strings。这与 GPU programming 中的 data residency 原则一致。加速效果最强的情况是：数据只复制到 device 一次，随后由多个 commands 处理，最后再 read back。

这个解释也说明了为什么 CUDA 并不总是在小输入上更快。即使 kernel 本身很高效，整个 command path 仍然要支付 launch、synchronization、output allocation 和 possible readback 的 fixed costs。因此，break-even point 是 system-level property，而不仅是 algorithmic property。

Concurrent geometry experiment 如图 \ref{fig:concurrency-throughput} 和图 \ref{fig:concurrency-latency} 所示。CPU throughput 随 worker 数量增加而提升，直到 contention 出现。CUDA throughput 更平缓，因为多个 script workers 会竞争同一个 GPU 和 synchronization point。

\begin{figure}[H]
  \centering
  \includegraphics[width=0.78\textwidth]{geometry-concurrency-throughput.png}
  \caption{并发几何 workload 的吞吐量（Throughput of concurrent geometry workloads）}
  \label{fig:concurrency-throughput}
\end{figure}

直接结论：增加 script workers 只能在 shared runtime 或 GPU resource 成为 bottleneck 之前提升 throughput。

\begin{figure}[H]
  \centering
  \includegraphics[width=0.78\textwidth]{geometry-concurrency-latency.png}
  \caption{并发几何 workload 的 P95 延迟（P95 latency）}
  \label{fig:concurrency-latency}
\end{figure}

直接结论：contention 会推高 P95 latency，因此后续需要 dedicated scheduling、memory pooling 和 stream-based CUDA execution。

\section{小结}
本章从 correctness、performance 和 scalability 三个维度评估 FTCL。Semantic regression 和 CPU/GPU equivalence tests 证明系统具有可验证的正确性基础；single-GPU 和 multi-GPU experiments 显示 CUDA backend 在大量 independent work items 上具有明显优势；ablation 和 concurrency experiments 则指出真正的瓶颈经常来自 allocation、readback、locking 和 scheduling。

The evaluation is deliberately end-to-end. The main lesson is that data residency, output size, and runtime synchronization dominate script-level heterogeneous performance.
